% Generated draft: requires mathematical and source review.
\documentclass[11pt]{article}
\usepackage[margin=25mm]{geometry}
\usepackage{fontspec}
\setmainfont{Latin Modern Roman}
\usepackage{amsmath,amssymb,mathtools,unicode-math}
\setmathfont{Latin Modern Math}
\usepackage{enumitem,needspace,xurl,hyperref}
\setlength{\parindent}{0pt}
\setlength{\parskip}{5pt}
\setlength{\emergencystretch}{3em}
\newcommand{\N}{\mathbb{N}}
\newcommand{\Z}{\mathbb{Z}}
\newcommand{\Q}{\mathbb{Q}}
\newcommand{\R}{\mathbb{R}}
\newcommand{\C}{\mathbb{C}}
\newcommand{\F}{\mathbb{F}}
\newcommand{\A}{\mathbb{A}}
\newcommand{\PP}{\mathbb{P}}
\newcommand{\E}{\mathbb{E}}
\newcommand{\eps}{\varepsilon}
\newcommand{\dd}{\,\mathrm{d}}
\begin{document}
\section*{UnsolvedMath ID: 701672}
\textbf{Source problem code:} catalogue entry
\subsection{Definitions and mathematical statement}
A variety of groups is a class of groups defined by a set of group identities, with group homomorphisms as morphisms. For groups \(A,B\in\mathfrak V\), their coproduct \(A*_{\mathfrak V}B\) is a group in \(\mathfrak V\) with homomorphisms \(\iota_A:A\to A*_{\mathfrak V}B\) and \(\iota_B:B\to A*_{\mathfrak V}B\) such that, for every \(H\in\mathfrak V\) and homomorphisms \(f:A\to H\), \(g:B\to H\), there is a unique homomorphism \(u:A*_{\mathfrak V}B\to H\) satisfying \(u\iota_A=f\) and \(u\iota_B=g\). An embedding means an injective group homomorphism. Let \(\operatorname{var}(G)\) denote the smallest variety containing \(G\), equivalently the variety defined by all group identities holding in \(G\). No finiteness assumptions are imposed.

\begin{enumerate}
\item[(a)] Do there exist a variety \(\mathfrak V\) and a group \(G\in\mathfrak V\) such that there is an embedding
\[
G*_{\mathfrak V}G\hookrightarrow G,
\]
but there is no embedding
\[
G*_{\mathfrak V}(G*_{\mathfrak V}G)\hookrightarrow G?
\]
Here the parenthesized group is, up to canonical isomorphism, the coproduct in \(\mathfrak V\) of three copies of \(G\).
\item[(b)] If the existence assertion in (a) holds, must there exist a group \(H\) for which
\[
H*_{\operatorname{var}(H)}H\hookrightarrow H
\]
and no embedding
\[
H*_{\operatorname{var}(H)}
  (H*_{\operatorname{var}(H)}H)\hookrightarrow H
\]
exists? The group \(H\) may differ from the group in an example for (a).
\end{enumerate}

The source emphasizes that an embedding \(e:G*_{\mathfrak V}G\hookrightarrow G\) does not automatically make the composite
\[
G*_{\mathfrak V}(G*_{\mathfrak V}G)
\xrightarrow{\,\mathrm{id}_G*_{\mathfrak V}e\,}
G*_{\mathfrak V}G
\xrightarrow{\,e\,}G
\]
an embedding: in a general group variety, the map induced by an embedding \(A\hookrightarrow B\) under \(G*_{\mathfrak V}(-)\) need not be injective. The nonexistence conditions above concern every possible embedding into the target group, not merely this particular composite. The registry labels the problem open; that label is supplied source information, not an independently verified assessment.
\subsection{Short English statement}
Can a group contain a copy of the coproduct of two copies of itself in some group variety, while containing no copy of the corresponding coproduct of three copies? If such an example exists, can one choose an example whose ambient variety is generated by the group itself?
\subsection{Sources}
\begin{itemize}
\item Local UnsolvedMath registry, permanent ID \(701672\), exact source code \texttt{catalogue entry}, title \emph{Kourovka Notebook Problem 20.14}. The supplied record attributes the problem to G. M. Bergman, with proposed year 2022, and labels its status \texttt{open}.
\item \emph{The Kourovka Notebook (Unsolved Problems in Group Theory)}, original problem 20.14; the supplied background identifies the 20th issue (2022) and the 21st issue edition (2026), arXiv:1401.0300v48. The supplied collection description names E. I. Khukhro and V. D. Mazurov as editors.
\item G. M. Bergman, \emph{Indag. Math.}, \textbf{18} (2007), 349--403. Supplied as a reference for part (a); no article title or URL is supplied.
\item G. M. Bergman, \emph{Algebra Number Theory}, \textbf{3} (2009), 847--879. Supplied as a reference for part (b) and related questions; no article title or URL is supplied.
\end{itemize}
\end{document}
